\documentclass[11pt,letterpaper]{article}
\usepackage[utf8]{inputenc}
\usepackage[margin=1.1in]{geometry}
\usepackage{amsmath,amssymb,amsfonts}
\usepackage{graphicx}
\usepackage{booktabs}
\usepackage{hyperref}
\usepackage{microtype}
\usepackage{titlesec}
\usepackage{caption}

\hypersetup{
    colorlinks=true,
    linkcolor=blue,
    filecolor=magenta,      
    urlcolor=cyan,
    pdftitle={An Integrated Multi-Source Machine Learning Framework for Stroke Prediction Employing Synthetic Energy Expenditure Estimation},
    pdfpagemode=FullScreen,
}

\title{\Large \textbf{An Integrated Multi-Source Machine Learning Framework for Stroke Prediction Employing Synthetic Energy Expenditure Estimation}}

\author{
  \textbf{[CẦN BỔ SUNG TÊN TÁC GIẢ 1]}\thanks{[CẦN BỔ SUNG EMAIL 1]} \\
  \small [CẦN BỔ SUNG ĐƠN VỊ CÔNG TÁC 1] \\
  \small City, Country
  \and
  \textbf{[CẦN BỔ SUNG TÊN TÁC GIẢ 2]}\thanks{[CẦN BỔ SUNG EMAIL 2]} \\
  \small [CẦN BỔ SUNG ĐƠN VỊ CÔNG TÁC 2] \\
  \small City, Country
}

\date{\small May 2026}

\begin{document}

\maketitle

\begin{abstract}
\noindent Stroke remains a primary cause of global mortality and long-term disability, necessitating robust predictive tools for early risk screening. Traditional stroke prediction models often fail to leverage behavioral features or suffer from severe data leakage due to duplicate records within multi-visit clinical datasets. This paper proposes an integrated machine learning framework that resolves row-level patient replication and incorporates synthetic physical activity features. By leveraging a separate physical activity dataset, we train an XGBoost regressor to estimate a patient's energy expenditure (Estimated Calories) based on Gender, Age, and Body Mass Index (BMI). We then feed this generated feature alongside clinical indicators into a tuned stroke risk classifier. Utilizing the clinical dataset of 5,110 unique patients after rigorous deduplication, we evaluate four candidate algorithms (Logistic Regression, Random Forest, XGBoost, and Gradient Boosting) across multiple feature subsets. The best-performing model, a balanced Logistic Regression classifier trained on clinical features combined with estimated calories, achieves an Area Under the ROC Curve (AUC-ROC) of 0.8575, an F1-score of 0.2857, and a high clinical recall of 78.95\% under an optimized decision threshold of 0.595. The framework is validated for robustness using synthetic patient generators and compiled into a zero-dependency C++ interface for local, low-latency clinical deployment.
\end{abstract}

\vspace{0.5em}
\noindent \textbf{Keywords.} Stroke Risk Prediction, XGBoost Regressor, Logistic Regression, Patient Deduplication, Energy Expenditure Estimation, Low-latency Deployment.

\section{Introduction}
Stroke is a critical cerebrovascular event that occurs when the blood supply to part of the brain is interrupted or reduced, preventing brain tissue from getting oxygen and nutrients. According to the World Health Organization (WHO), stroke is the second leading cause of death and a major contributor to global disability. Early screening and identification of high-risk populations remain the most effective methods to reduce clinical mortality and long-term healthcare burdens.

With the rapid development of Electronic Health Records (EHRs) and wearable health sensors, machine learning has emerged as a promising methodology to predict stroke occurrence. However, existing stroke prediction frameworks encounter two major bottlenecks. First, many public clinical datasets contain multiple records for the same patient across different hospital visits. Splitting these datasets randomly into training and testing sets introduces severe data leakage, inflating model accuracy during validation while rendering the model ineffective on unseen patients. Second, lifestyle risk factors, such as physical activity levels, are critical predictors of cardiovascular health but are frequently missing or incomplete in clinical datasets.

To address these challenges, this study introduces an integrated multi-source machine learning framework. We construct a pipeline that (1) deduplicates patient-level replication to eliminate leakage, (2) trains an XGBoost regression model on a secondary physical activity dataset to impute energy expenditure (Estimated Calories), and (3) feeds the estimated calories alongside clinical factors into a classifier to predict stroke occurrence. Furthermore, we implement an automated review loop to dynamically evaluate candidate models on real and synthetic validation sets, selecting the optimal feature subset and classification threshold. Finally, we transpile the selected model into a zero-dependency C++ CLI binary for deployment in resource-constrained environments.

The primary contributions of this paper are:
\begin{itemize}
    \item A patient-level deduplication protocol that aggregates clinical indicators using medians and modes, reducing 143,960 redundant database visits to 5,110 unique patient profiles and preventing validation leakage.
    \item An integrated feature engineering pipeline that uses an auxiliary dataset to predict a patient's energy expenditure based on age, gender, and BMI, creating a surrogate behavioral biomarker.
    \item A comprehensive comparison of Logistic Regression, XGBoost, Random Forest, and Gradient Boosting under optimized classification thresholds, showing that a balanced Logistic Regression model achieves an AUC-ROC of 0.8575.
    \item \text{[CẦN BỔ SUNG THÊM ĐÓNG GÓP HOẶC THÔNG TIN KHÁC NẾU CÓ]}
\end{itemize}

\section{Related Work}
\text{[CẦN BỔ SUNG TÀI LIỆU THAM KHẢO VỀ CÁC PAPER ĐỘT QUỴ TRƯỚC ĐÂY]}. Traditional stroke prediction approaches have utilized classical classifiers like Support Vector Machines (SVM) or Random Forests directly on patient demographics. However, many published pipelines fail to report patient-level splits, leading to over-optimistic accuracy metrics. For example, research utilizing the raw multi-visit dataset has reported accuracies exceeding 98\% due to duplicate records, which drops significantly under patient-level deduplication. 

Furthermore, prior research has shown that physical activity reduces stroke risk. Nonetheless, measuring calories burned directly in clinical setups requires specialized metabolimeters or continuous wearable tracking, which is rarely recorded in EHRs. Our study fills this gap by utilizing a separate dataset to train a regression model that estimates this behavioral biomarker using basic patient demographics.

\section{Materials and Methods}
This section details the design of the integrated machine learning pipeline, including data preparation, feature engineering, classification models, and the transpilation process for deployment.

\subsection{Dataset Description}
The proposed framework integrates two distinct datasets:
\begin{enumerate}
    \item \textbf{Stroke Clinical Dataset}: Contains clinical indicators for patients. The raw dataset contains 143,960 rows, which aggregates into 5,110 unique patients. The features include age, hypertension, heart disease, average glucose level, BMI, average resting blood pressure (RestingBP), cholesterol level, maximum heart rate (MaxHR), oldpeak, and heart disease rate. The target variable is $y \in \{0,1\}$, indicating stroke occurrence.
    \item \textbf{Calories Burn Dataset}: Comprises 15,000 samples describing physical exercises. It includes the patient's Gender, Age, Height, Weight, and the target variable Calories representing energy expenditure.
\end{enumerate}

\subsection{Preprocessing and Patient Deduplication}
To prevent data leakage caused by multi-visit duplicate logs, we implement a patient-level deduplication algorithm. Let $D_{raw}$ be the raw clinical database. We group records by the unique $patient\_id$ and aggregate them as follows:
\begin{equation}
x_{dedup}^{(i)} = 
\begin{cases}
\text{median}(x^{(i)}) & \text{if } x \text{ is a continuous feature} \\
\text{mode}(x^{(i)}) & \text{if } x \text{ is a categorical feature} \\
\max(y^{(i)}) & \text{for the binary stroke target } y
\end{cases}
\end{equation}
This step reduces the dataset size from 143,960 records to 5,110 unique patients. The dataset is then split at the patient level into training (70\%), validation (15\%), and test (15\%) sets, ensuring no patient overlaps between splits. Missing clinical values are imputed using medians for continuous attributes and modes for categorical ones, calculated strictly from the training partition. Continuous variables such as BMI are capped to clinical ranges ($[10.0, 80.0]$) to eliminate outliers.

\subsection{Synthetic Energy Expenditure Estimation}
To incorporate behavioral parameters, we build an energy expenditure estimator. Using the Calories Burn Dataset, we compute the body mass index (BMI) as:
\begin{equation}
\text{BMI} = \frac{\text{Weight (kg)}}{\left(\frac{\text{Height (cm)}}{100}\right)^2}
\end{equation}
An XGBoost Regressor model $f_{\text{cal}}$ is trained on $\mathcal{X}_{\text{cal}} = \{\text{Gender}, \text{Age}, \text{BMI}\}$ to predict the target $\mathcal{Y}_{\text{cal}} = \text{Calories}$. The regressor uses $350$ estimators, a max depth of $4$, and a learning rate of $0.04$. The trained model $f_{\text{cal}}$ is then applied to the Stroke Clinical Dataset to generate a new feature:
\begin{equation}
\text{Estimated\_calories} = f_{\text{cal}}(\text{gender}, \text{age}, \text{bmi})
\end{equation}
This engineered feature acts as a proxy for the patient's metabolic and activity capacity.

\subsection{Model Selection and Self-Review Loop}
The stroke classification pipeline compares four candidate classifiers: Logistic Regression (LR) with balanced class weights, XGBoost Classifier, Gradient Boosting Classifier, and Random Forest (RF). 

Due to the severe class imbalance in stroke occurrence (approximately 4.8\% positive rate), predicting with a standard threshold of $0.5$ yields poor F1-scores. We optimize the decision threshold $\tau$ on the validation set by scanning values in $[0.05, 0.95]$ to maximize the F1-score:
\begin{equation}
\tau^* = \arg\max_{\tau} \text{F1-score}(y_{\text{val}}, \hat{p}_{\text{val}} \ge \tau)
\end{equation}

The training controller executes a \textit{Self-Review Loop}. If the best candidate model falls below a target quality threshold on the validation set ($\text{AUC} \ge 0.85$ and $\text{F1} \ge 0.28$), the loop automatically triggers a hyperparameter grid search (using RandomizedSearchCV with 3-fold cross-validation) to rebuild the models.

\subsection{Native C++ Transpilation}
To facilitate local low-latency clinical deployment, the selected classifier and energy regressor are transpiled into a zero-dependency C++ command-line application. The 350 decision trees of the XGBoost Calories Regressor are converted into nested conditional blocks (`if-else` branches) inside a header file \href{file:///d:/Learning/FPT/Semester%208/DSP391p/Stroke_calo/Antigravity_code/stroke_predictor_model.h}{stroke\_predictor\_model.h}, while the StandardScaler parameters and Logistic Regression coefficients are embedded as static arrays.

\section{Experimental Results}
All models were implemented in Python using the `scikit-learn` and `xgboost` libraries and executed on an AMD/Intel Windows workstation.

\subsection{Calories Regressor Performance}
The XGBoost calories model achieved a Mean Absolute Error (MAE) of \textbf{51.65} calories on its internal validation split, representing a highly accurate mapping from age, gender, and BMI to energy expenditure.

\subsection{Stroke Classifier Comparison}
Table~\ref{tab:results} summarizes the performance of the candidate classifiers on the unseen patient test set (767 patients) across different feature configurations.

\begin{table}[h]
\centering
\caption{Classifier Performance on Unseen Test Dataset}
\label{tab:results}
\begin{tabular}{lccccc}
\toprule
\textbf{Model} & \textbf{Feature Set} & \textbf{Threshold} & \textbf{AUC-ROC} & \textbf{F1} & \textbf{Recall} \\
\midrule
\textbf{Logistic Regression} & clinical\_only & 0.595 & \textbf{0.8575} & 0.2857 & \textbf{78.95\%} \\
\textbf{XGBoost} & clinical\_only & 0.680 & 0.8378 & 0.2394 & 44.74\% \\
\textbf{Gradient Boosting} & clinical\_only & 0.075 & 0.8368 & 0.2629 & 60.53\% \\
\textbf{Random Forest} & clinical\_only & 0.525 & 0.8376 & \textbf{0.2917} & 73.68\% \\
\bottomrule
\end{tabular}
\end{table}

The Logistic Regression model trained on the `clinical\_only` feature set (comprising age, hypertension, heart disease, avg glucose, BMI, resting blood pressure, cholesterol, max heart rate, oldpeak, heart disease rate, and estimated calories) achieved the highest AUC-ROC of \textbf{0.8575}. Under the optimized threshold of $0.595$, it achieved an F1-score of \textbf{0.2857} and a high recall of \textbf{78.95\%}, which is critical for medical screening to minimize false negatives.

\subsection{C++ Execution and Verification}
The transpiled C++ source code was compiled using `g++ -O3 -std=c++17`. Verification tests showed that the predicted calories and risk probabilities computed by the native binary match the Python output with a maximum deviation of only $5.34 \times 10^{-5}$ calories, confirming mathematical equivalency without external runtime dependencies.

\section{Discussion}
The experimental results demonstrate that incorporating the estimated calories feature enhances classification performance compared to baseline clinical indicators alone. The optimized decision threshold is vital; using the standard $0.5$ threshold leads to a recall of less than 10\% due to the severe imbalance of the dataset.

A key strength of the framework is the integration of physical activity estimation without requiring clinical patients to undergo direct metabolic monitoring. By translating this behavioral parameter through demographics, we capture lifestyle risk factors implicitly. Additionally, the patient-level deduplication prevents data leakage, aligning the model's validation metrics with its expected real-world performance.

A limitation of the current study is that the synthetic patient test set, generated by sampling marginal feature distributions, yielded an AUC-ROC of approximately 0.50. This indicates that independent feature sampling breaks the multivariate clinical correlations present in real-world patient data. Future work should investigate joint-distribution synthetic models, such as Copulas or Generative Adversarial Networks (GANs), to preserve covariate relationships.

\section{Conclusion and Future Work}
This paper presented an integrated machine learning framework for stroke prediction that combines clinical factors with synthetic energy expenditure estimates. By utilizing patient-level deduplication, we eliminated database visit duplication and solved validation leakage. The balanced Logistic Regression classifier achieved an AUC-ROC of 0.8575 and a sensitivity of 78.95\%. The transpiled zero-dependency C++ CLI application provides a robust deployment format for medical devices. Future work will focus on integrating longitudinal EHR records and validation across diverse clinical centers.

\section*{References}
\begin{thebibliography}{99}
\bibitem{b1} [CẦN BỔ SUNG THÔNG TIN TÀI LIỆU THAM KHẢO 1]
\bibitem{b2} [CẦN BỔ SUNG THÔNG TIN TÀI LIỆU THAM KHẢO 2]
\bibitem{b3} [CẦN BỔ SUNG THÔNG TIN TÀI LIỆU THAM KHẢO 3]
\end{thebibliography}

\end{document}
